\documentclass[aspectratio=169,11pt]{beamer}
\usepackage{../../../shared/theme/esmad_beamer_theme}
\usepackage{amsmath,amssymb}
\usepackage{booktabs}

\title[Dynamic Treatment Effects]{Dynamic Treatment Effects}
\subtitle{Panel and Time-Series Causal Responses}
\date{\today}

\newcommand{\E}{\mathbb{E}}

\begin{document}

\begin{frame}
\titlepage
\end{frame}

\begin{frame}{Learning Goals}
\begin{itemize}
\item distinguish immediate, delayed, cumulative, and long-run treatment effects;
\item formulate distributed-lag and event-time causal estimands;
\item explain anticipation, carry-over, and persistence;
\item recognize dynamic confounding and treatment feedback;
\item understand when lagged outcomes create additional identification problems;
\item summarize treatment-response paths without collapsing them into misleading static coefficients.
\end{itemize}
\end{frame}

\begin{frame}{Outline}
\tableofcontents
\end{frame}

\section{Why Effects Are Dynamic}

\begin{frame}{A Static Effect Can Hide the Mechanism}
A static model writes
\[
Y_{it}=\alpha_i+\lambda_t+\beta D_{it}+u_{it}.
\]

But treatment may act through several future periods.

A richer representation is
\[
Y_{it}
=
\alpha_i+\lambda_t
+
\sum_{k=0}^{K}\beta_kD_{i,t-k}
+
u_{it}.
\]

\begin{alertblock}{Key question}
Which horizon is scientifically meaningful: now, next period, or the cumulative effect over several periods?
\end{alertblock}
\end{frame}

\begin{frame}{Types of Dynamic Responses}
Treatment effects may show:
\begin{itemize}
\item immediate impact;
\item delayed onset;
\item gradual accumulation;
\item temporary spike;
\item fade-out;
\item persistent long-run shift.
\end{itemize}

Different mechanisms imply different estimands and different modeling choices.
\end{frame}

\section{Distributed Lags}

\begin{frame}{Distributed-Lag Model}
Consider
\[
Y_t
=
\alpha
+
\sum_{k=0}^{K}\beta_k D_{t-k}
+
u_t.
\]

Each \(\beta_k\) captures the marginal effect associated with treatment \(k\) periods earlier, conditional on the specification.

\begin{block}{Immediate effect}
\[
\beta_0.
\]
\end{block}
\end{frame}

\begin{frame}{Cumulative Effects}
The cumulative effect through horizon \(h\) is
\[
\Theta_h
=
\sum_{k=0}^{h}\beta_k.
\]

The long-run finite-horizon effect is
\[
\Theta_K
=
\sum_{k=0}^{K}\beta_k.
\]

\begin{alertblock}{Interpretation}
A small immediate coefficient can coexist with a large cumulative effect.
\end{alertblock}
\end{frame}

\begin{frame}{Choosing the Lag Length}
Too few lags:
\begin{itemize}
\item omit delayed effects;
\item bias cumulative responses;
\item force persistence into the disturbance.
\end{itemize}

Too many lags:
\begin{itemize}
\item inflate variance;
\item create multicollinearity;
\item make individual coefficients unstable.
\end{itemize}
\end{frame}

\section{Event-Time Effects}

\begin{frame}{Dynamic Event-Time Estimands}
For treatment cohort \(g\), define
\[
ATT(g,t)
=
\E[Y_t(1)-Y_t(0)\mid T_i=g].
\]

Relative event time is
\[
k=t-g.
\]

Then dynamic effects can be summarized as functions of \(k\), rather than calendar time alone.
\end{frame}

\begin{frame}{Why Event Time Helps}
Event-time effects answer questions such as:
\begin{itemize}
\item Does the effect begin immediately?
\item How quickly does it grow?
\item Does it persist?
\item When does it plateau?
\item Does it reverse?
\end{itemize}

This is often more useful for policy or operational decisions than a single post-treatment average.
\end{frame}

\section{Anticipation}

\begin{frame}{Treatment Can Start Before the Official Date}
Suppose treatment is announced at \(T_a\) and implemented at \(T_0>T_a\).

If units react after announcement,
\[
Y_t(1)\neq Y_t(0)
\]
for some \(t<T_0\).

\begin{alertblock}{Consequence}
Nominal pre-treatment observations may already contain causal effects.
\end{alertblock}
\end{frame}

\begin{frame}{Anticipation Examples}
\begin{itemize}
\item customers buy early before a tax increase;
\item firms hire before a subsidy begins;
\item investors reprice assets after an announcement;
\item hospitals change scheduling before a new rule is enforced.
\end{itemize}

The relevant treatment timing is behavioral, not merely administrative.
\end{frame}

\section{Carry-Over and Persistence}

\begin{frame}{Treatment Can Persist After Exposure Ends}
If treatment ends at \(T_1\), the effect may continue:
\[
Y_t(1)\neq Y_t(0)
\qquad
t>T_1.
\]

This can occur through:
\begin{itemize}
\item habit formation;
\item capital accumulation;
\item learning;
\item biological persistence;
\item institutional inertia.
\end{itemize}
\end{frame}

\begin{frame}{Washout Periods}
When carry-over is plausible, one may need:
\begin{itemize}
\item explicit lagged treatment terms;
\item post-treatment washout windows;
\item separate estimands for exposure and persistence;
\item sensitivity to assumed decay patterns.
\end{itemize}
\end{frame}

\section{Dynamic Confounding}

\begin{frame}{Past Outcomes Can Affect Future Treatment}
Suppose
\[
D_{it}
=
f(Y_{i,t-1},X_{it},\eta_{it}).
\]

Then past outcomes affect treatment assignment.

If past outcomes are themselves affected by earlier treatment, standard regression adjustment can become problematic.

\begin{alertblock}{Feedback}
Treatment changes outcomes, outcomes influence later treatment, and treatment history becomes endogenous.
\end{alertblock}
\end{frame}

\begin{frame}{Time-Varying Confounders Affected by Prior Treatment}
Let \(L_t\) be a time-varying covariate such that
\[
D_{t-1}\rightarrow L_t,
\qquad
L_t\rightarrow D_t,
\qquad
L_t\rightarrow Y.
\]

Adjusting for \(L_t\) can block part of the treatment effect while failing to adjust correctly for dynamic confounding.

This motivates longitudinal methods such as marginal structural models in appropriate settings.
\end{frame}

\section{Lagged Outcomes}

\begin{frame}{Dynamic Panel Model}
Consider
\[
Y_{it}
=
\rho Y_{i,t-1}
+
\beta D_{it}
+
\alpha_i
+
u_{it}.
\]

Lagged outcomes can absorb persistence, but the within transformation creates correlation between transformed lagged outcomes and transformed errors.
\end{frame}

\begin{frame}{Nickell Bias}
In short panels, fixed-effects estimation of
\[
Y_{it}=\rho Y_{i,t-1}+\alpha_i+u_{it}
\]
is biased because
\[
Y_{i,t-1}-\bar Y_i
\]
correlates with
\[
u_{it}-\bar u_i.
\]

The bias decreases as \(T\) grows, but can be substantial for small \(T\).
\end{frame}

\section{Impulse-Response Style Thinking}

\begin{frame}{Treatment-Response Path}
A dynamic causal response can be summarized as
\[
\{\tau_0,\tau_1,\tau_2,\ldots,\tau_H\}.
\]

This resembles an impulse response:
\begin{itemize}
\item \(\tau_0\): immediate effect;
\item \(\tau_h\): effect \(h\) periods later;
\item cumulative response: \(\sum_{j=0}^{h}\tau_j\).
\end{itemize}
\end{frame}

\begin{frame}{Do Not Confuse Forecasting IRFs with Causal IRFs}
A dynamic response is causal only if the shock or treatment variation is identified causally.

A predictive time-series model can describe temporal propagation without establishing exogenous intervention effects.

\begin{alertblock}{Identification first}
Dynamic machinery does not substitute for a causal design.
\end{alertblock}
\end{frame}

\section{Heterogeneous Dynamics}

\begin{frame}{Different Units May Respond Differently}
Let
\[
\tau_{i,h}
=
Y_{i,t+h}(1)-Y_{i,t+h}(0).
\]

Heterogeneity may occur across:
\begin{itemize}
\item regions;
\item firms;
\item age groups;
\item treatment intensity;
\item baseline risk;
\item adoption cohorts.
\end{itemize}
\end{frame}

\begin{frame}{Average Dynamic Effects}
An average event-time effect is
\[
\tau_h
=
\E[\tau_{i,h}].
\]

But the composition of units contributing to \(\tau_h\) can change with horizon.

\begin{alertblock}{Comparability issue}
At long horizons, only early-treated cohorts may remain observable.
\end{alertblock}
\end{frame}

\section{Applied Reasoning}

\begin{frame}{Example: Regional Pricing Policy}
Suppose a regional pricing intervention starts in January.

Possible effect path:
\begin{itemize}
\item month 0: little response;
\item months 1--2: customers begin adapting;
\item months 3--6: demand stabilizes at a new level;
\item later months: competitor responses partially offset the effect.
\end{itemize}

A single post-treatment coefficient would compress all of this into one number.
\end{frame}

\begin{frame}{A Credible Dynamic-Effects Workflow}
\begin{enumerate}
\item Define the treatment start, announcement date, and exposure duration.
\item Specify the horizons of scientific interest.
\item State immediate and cumulative estimands.
\item Plot event-time responses.
\item Include lags consistent with the mechanism.
\item Test sensitivity to lag length.
\item Examine anticipation and carry-over.
\item Address dynamic confounding and feedback.
\item Report heterogeneity across cohorts or subgroups.
\end{enumerate}
\end{frame}

\begin{frame}{Common Failure Modes}
\begin{itemize}
\item forcing a static model onto a delayed-response process;
\item choosing lag length only by significance;
\item ignoring anticipation;
\item treating persistent post-exposure effects as new treatments;
\item controlling naively for time-varying mediators/confounders;
\item interpreting lagged-outcome FE estimates without considering dynamic-panel bias;
\item reporting only pointwise effects without cumulative interpretation.
\end{itemize}
\end{frame}

\begin{frame}{Synthesis: Time Is Part of the Estimand}
\begin{itemize}
\item Treatment effects can unfold over several periods.
\item Distributed-lag models decompose immediate and delayed responses.
\item Cumulative effects often matter more than a single-period coefficient.
\item Event-time estimands reveal onset, persistence, and decay.
\item Anticipation shifts the effective treatment date earlier.
\item Carry-over extends effects beyond the exposure period.
\item Dynamic confounding and lagged outcomes can invalidate naive static adjustment.
\item Causal identification must remain valid at every horizon being interpreted.
\end{itemize}
\end{frame}

\begin{frame}{Selected References}
\small
\begin{itemize}
\item Arellano, M. (2003). \textit{Panel Data Econometrics}. Oxford University Press.
\item Nickell, S. (1981). Biases in Dynamic Models with Fixed Effects. \textit{Econometrica}, 49(6), 1417--1426.
\item Robins, J. M., Hernan, M. A., and Brumback, B. (2000). Marginal Structural Models and Causal Inference in Epidemiology. \textit{Epidemiology}, 11(5), 550--560.
\item Sun, L., and Abraham, S. (2021). Estimating Dynamic Treatment Effects in Event Studies with Heterogeneous Treatment Effects. \textit{Journal of Econometrics}, 225(2), 175--199.
\item Callaway, B., and Sant'Anna, P. H. C. (2021). Difference-in-Differences with Multiple Time Periods. \textit{Journal of Econometrics}, 225(2), 200--230.
\item Jordà, O. (2005). Estimation and Inference of Impulse Responses by Local Projections. \textit{American Economic Review}, 95(1), 161--182.
\end{itemize}
\end{frame}

\contactslide

\end{document}
